\frontsection{Danh mục thuật ngữ}

\begin{spacing}{1}
\small
\renewcommand{\arraystretch}{1.12}
\begin{longtable}{@{}>{\bfseries\color{HCMUTDarkBlue}}p{4.1cm}p{10.7cm}@{}}
  \toprule
  \multicolumn{1}{@{}l}{\textbf{\color{HCMUTDarkBlue}Thuật ngữ}} &
  \textbf{Giải thích trong phạm vi báo cáo} \\
  \midrule
  \endfirsthead

  \toprule
  \multicolumn{1}{@{}l}{\textbf{\color{HCMUTDarkBlue}Thuật ngữ}} &
  \textbf{Giải thích trong phạm vi báo cáo} \\
  \midrule
  \endhead

  \bottomrule
  \endfoot

  Activity
    & Thành phần Android đại diện cho một màn hình hoặc điểm tương tác chính với người dùng và có vòng đời do hệ điều hành quản lý. \\
  Actor
    & Tác nhân tương tác với hệ thống trong mô hình use case, có thể là người sử dụng hoặc một hệ thống bên ngoài. \\
  Annotation
    & Thông tin chú thích gắn với dữ liệu, chẳng hạn nhãn lớp, bounding box, keypoint hoặc thời điểm xảy ra hành động trong video. \\
  Artifact
    & Sản phẩm có thể lưu trữ và kiểm tra được tạo ra trong quá trình phát triển, chẳng hạn tệp mô hình, ứng dụng cài đặt, nhật ký hoặc kết quả kiểm thử. \\
  Augmentation
    & Kỹ thuật tạo thêm biến thể từ dữ liệu huấn luyện bằng các phép biến đổi như lật ảnh, thay đổi độ sáng, co giãn hoặc xoay ảnh. \\
  Backbone
    & Phần mạng nơ-ron đảm nhiệm trích xuất đặc trưng nền tảng từ dữ liệu đầu vào trước khi các đầu dự đoán thực hiện nhiệm vụ cụ thể. \\
  Batch
    & Nhóm mẫu dữ liệu được đưa qua mô hình trong cùng một bước huấn luyện hoặc suy luận. \\
  Bonding
    & Cơ chế lưu thông tin khóa sau khi hai thiết bị BLE ghép đôi để có thể tái lập liên kết bảo mật ở những lần kết nối sau. \\
  Bounding box
    & Hình chữ nhật bao quanh đối tượng trong ảnh, dùng để biểu diễn vị trí và kích thước của đối tượng được phát hiện. \\
  Callback
    & Hàm được đăng ký để hệ thống gọi lại khi một tác vụ hoàn tất hoặc khi một sự kiện tương ứng xảy ra. \\
  Calibration
    & Quá trình hiệu chỉnh tham số, ngưỡng hoặc kết quả đo để hệ thống phù hợp hơn với thiết bị và điều kiện vận hành cụ thể. \\
  Candidate
    & Địa chỉ IP và cổng có khả năng tạo đường truyền giữa hai đầu cuối trong quá trình thiết lập kết nối WebRTC bằng ICE. \\
  Characteristic
    & Thành phần dữ liệu nằm trong một GATT service, có giá trị, định danh và các thuộc tính quy định thao tác đọc, ghi hoặc thông báo. \\
  Closed-loop control
    & Điều khiển vòng kín; đầu ra thực tế được đo và phản hồi liên tục để bộ điều khiển điều chỉnh tín hiệu tác động theo giá trị đặt. \\
  Codec
    & Phương thức mã hóa và giải mã âm thanh hoặc hình ảnh để phục vụ lưu trữ và truyền dữ liệu media. \\
  Confidence score
    & Giá trị thể hiện mức độ tin cậy của mô hình đối với một kết quả dự đoán, thường được dùng để lọc các phát hiện không chắc chắn. \\
  Data channel
    & Kênh truyền dữ liệu tùy ý giữa hai đầu cuối WebRTC, hoạt động bên cạnh các luồng âm thanh và hình ảnh. \\
  Data split
    & Cách phân chia dữ liệu thành các tập huấn luyện, xác thực và kiểm tra nhằm huấn luyện mô hình và đánh giá khả năng khái quát. \\
  Dataset
    & Tập dữ liệu được tổ chức để huấn luyện, xác thực hoặc đánh giá thuật toán và mô hình học máy. \\
  Delegate
    & Thành phần chuyển một phần hoặc toàn bộ đồ thị suy luận sang bộ tăng tốc như GPU hoặc NPU thay vì chỉ thực thi trên CPU. \\
  Descriptor
    & Thành phần mô tả bổ sung cho characteristic trong GATT, cung cấp thông tin cấu hình hoặc cách diễn giải giá trị. \\
  Differential drive
    & Cơ cấu dẫn động vi sai sử dụng hai bánh chủ động độc lập; hướng và vận tốc robot được xác định từ chênh lệch vận tốc giữa hai bánh. \\
  Drop frame
    & Việc chủ động bỏ qua một khung hình khi hệ thống chưa xử lý xong dữ liệu trước đó, nhằm tránh hàng đợi và độ trễ tiếp tục tăng. \\
  Edge AI
    & Cách triển khai trong đó mô hình AI được thực thi gần nguồn tạo dữ liệu, chẳng hạn trực tiếp trên điện thoại gắn trên robot. \\
  Embedding
    & Vector đặc trưng số biểu diễn nội dung hoặc đặc điểm của dữ liệu, có thể được dùng để so sánh mức độ tương đồng giữa các đối tượng. \\
  End-to-end
    & Cách tổ chức mô hình hoặc quy trình xử lý trực tiếp từ đầu vào đến kết quả cuối với ít khâu trung gian được xây dựng thủ công. \\
  Epoch
    & Một lượt mô hình sử dụng toàn bộ tập dữ liệu huấn luyện để cập nhật các tham số. \\
  Feature extraction
    & Quá trình trích xuất các đặc trưng có ý nghĩa từ dữ liệu đầu vào để phục vụ nhận dạng, phân loại hoặc theo dõi. \\
  Fine-tuning
    & Quá trình tiếp tục huấn luyện một mô hình đã được huấn luyện trước trên dữ liệu của bài toán cụ thể. \\
  Firmware
    & Chương trình chạy trên vi điều khiển, đảm nhiệm giao tiếp với phần cứng và thực hiện các chức năng điều khiển ở mức thiết bị. \\
  Foreground service
    & Dịch vụ Android thực hiện tác vụ người dùng có thể nhận biết trong thời gian dài và phải duy trì thông báo thường trực. \\
  Frame
    & Một khung hình đơn trong chuỗi video; nhiều frame liên tiếp tạo thành dữ liệu chuyển động theo thời gian. \\
  Fragment
    & Thành phần giao diện Android có vòng đời riêng và thường được đặt bên trong Activity để tổ chức từng khu vực chức năng của màn hình. \\
  GATT service
    & Nhóm các characteristic phục vụ cùng một chức năng trong mô hình dữ liệu GATT và được nhận diện bằng UUID. \\
  Ground truth
    & Nhãn tham chiếu được xem là đúng, dùng làm cơ sở huấn luyện hoặc đối chiếu khi đánh giá kết quả dự đoán. \\
  Heuristic
    & Quy tắc kinh nghiệm giúp đưa ra quyết định nhanh dựa trên đặc điểm của bài toán nhưng không bảo đảm tối ưu trong mọi trường hợp. \\
  Hungarian algorithm
    & Thuật toán tìm phương án gán có tổng chi phí nhỏ nhất, được dùng để liên kết kết quả phát hiện với các đối tượng đang được theo dõi. \\
  Hysteresis
    & Cơ chế sử dụng ngưỡng chuyển vào và ngưỡng chuyển ra khác nhau để hạn chế trạng thái thay đổi liên tục khi giá trị dao động gần biên. \\
  Inference
    & Quá trình sử dụng mô hình đã huấn luyện để tạo kết quả dự đoán từ dữ liệu đầu vào mới. \\
  Interpreter
    & Thành phần nạp mô hình, quản lý tensor và thực thi các toán tử của mô hình trong môi trường suy luận. \\
  Kalman filter
    & Bộ lọc ước lượng trạng thái bằng cách kết hợp dự báo từ mô hình chuyển động với thông tin quan sát mới. \\
  Keypoint
    & Điểm mốc biểu diễn vị trí một bộ phận trên cơ thể người, chẳng hạn vai, hông, gối hoặc mắt cá chân. \\
  Letterbox
    & Phép đổi kích thước ảnh nhưng vẫn giữ tỷ lệ khung hình bằng cách bổ sung vùng đệm vào phần còn thiếu. \\
  Listener
    & Thành phần đăng ký để nhận thông báo khi dữ liệu hoặc trạng thái được theo dõi có thay đổi. \\
  Media stream
    & Luồng dữ liệu âm thanh hoặc hình ảnh được thu nhận và truyền liên tục trong một phiên liên lạc thời gian thực. \\
  Metadata
    & Dữ liệu mô tả bổ sung về cấu trúc, định dạng hoặc ý nghĩa của một tệp, mô hình hay tập dữ liệu. \\
  Model quantization
    & Kỹ thuật giảm độ chính xác biểu diễn của trọng số hoặc tensor nhằm giảm dung lượng mô hình và chi phí tính toán. \\
  Object detection
    & Bài toán xác định lớp và vị trí của một hoặc nhiều đối tượng trong ảnh hoặc video. \\
  Object tracking
    & Bài toán duy trì định danh và ước lượng vị trí của đối tượng qua nhiều khung hình liên tiếp. \\
  Offer--Answer
    & Cơ chế thương lượng phiên trong WebRTC, trong đó một đầu cuối gửi đề nghị cấu hình và đầu cuối còn lại phản hồi bằng cấu hình tương thích. \\
  On-device AI
    & Hình thức thực thi mô hình AI trực tiếp trên thiết bị đầu cuối thay vì gửi dữ liệu đến máy chủ để suy luận. \\
  Payload
    & Phần dữ liệu nội dung được mang trong một bản tin truyền thông, không bao gồm các thông tin điều khiển của giao thức. \\
  Pipeline
    & Chuỗi các bước xử lý có thứ tự, trong đó đầu ra của một bước trở thành đầu vào cho bước tiếp theo. \\
  Pose estimation
    & Bài toán ước lượng tư thế cơ thể thông qua vị trí và độ tin cậy của các keypoint trên ảnh. \\
  Prompt
    & Nội dung hướng dẫn hoặc ngữ cảnh được gửi đến mô hình ngôn ngữ để định hướng cách tạo phản hồi. \\
  Preprocessing/Post-processing
    & Tiền xử lý chuẩn bị dữ liệu trước khi đưa vào mô hình; hậu xử lý diễn giải và lọc kết quả sau suy luận. \\
  Re-identification
    & Quá trình nhận dạng lại một đối tượng đã xuất hiện trước đó sau khi bị che khuất, mất dấu hoặc rời khỏi vùng quan sát. \\
  Setpoint
    & Giá trị đặt mà hệ thống điều khiển hướng tới, chẳng hạn vận tốc mong muốn của từng bánh xe. \\
  Signaling
    & Quá trình trao đổi thông tin điều khiển để thiết lập và duy trì phiên WebRTC; dữ liệu media không truyền qua kênh này sau khi phiên được hình thành. \\
  Softmax
    & Hàm chuyển các giá trị đầu ra của mô hình thành phân bố xác suất trên các lớp dự đoán. \\
  Streaming
    & Phương thức truyền và xử lý dữ liệu liên tục theo từng phần thay vì chờ toàn bộ nội dung hoàn tất. \\
  Tensor
    & Cấu trúc dữ liệu nhiều chiều dùng để biểu diễn đầu vào, đầu ra và các giá trị trung gian của mô hình học sâu. \\
  Threshold
    & Giá trị ngưỡng dùng để quyết định giữ, loại hoặc chuyển trạng thái dựa trên một đại lượng đo hay kết quả dự đoán. \\
  Timestamp
    & Giá trị biểu diễn thời điểm một dữ liệu hoặc sự kiện được tạo ra, dùng để sắp xếp và kiểm tra tính còn hiệu lực. \\
  Track/tracklet
    & Chuỗi các quan sát được liên kết và gán cùng một định danh; tracklet thường chỉ một đoạn theo dõi ngắn hoặc chưa hoàn chỉnh. \\
  Tracking-by-detection
    & Hướng theo dõi sử dụng kết quả phát hiện ở từng khung hình, sau đó dự báo và liên kết các kết quả để duy trì định danh theo thời gian. \\
  Transfer learning
    & Kỹ thuật tận dụng tri thức từ mô hình đã huấn luyện trên một bài toán hoặc tập dữ liệu khác để huấn luyện bài toán mới. \\
  Use case
    & Mô tả một mục tiêu tương tác giữa tác nhân và hệ thống, bao gồm các điều kiện và luồng hành vi cần thiết để đạt mục tiêu đó. \\
  Watchdog
    & Cơ chế giám sát thời gian cập nhật; khi lệnh hoặc tín hiệu không được làm mới trong khoảng quy định, hệ thống chuyển về trạng thái an toàn. \\
  Write without response
    & Phương thức ghi GATT không yêu cầu phản hồi xác nhận ở tầng ATT, phù hợp với dữ liệu được cập nhật thường xuyên và cần độ trễ thấp. \\
\end{longtable}
\end{spacing}
